\documentclass[dvipdfmx,a4paper]{jsarticle}
\usepackage{amsmath,amssymb,amsfonts,amsthm}
\usepackage[dvipdfmx]{graphicx}
\usepackage{tikz}
\usetikzlibrary{positioning, intersections, calc, arrows.meta, math, through, shadows,automata}
\usepackage{caption}
\usepackage{tcolorbox}
\tcbuselibrary{theorems,breakable}
\usepackage{enumerate}
\usepackage{mathtools}
\usepackage{otf}
\usepackage{xspace}
\usepackage{newpxtext}
\usepackage[utf8]{inputenc}
\usepackage{CJKutf8,CJKspace,CJKpunct}
\usepackage{pgfplots}
\usepackage{lastpage}
\usepackage{fancyhdr}
\pagestyle{fancy}
\fancyhf{}
\fancyhead[L]{\textsc{特殊講義（専門01)}}
\fancyhead[R]{\texttt{YI Ran} - $\mathnormal{61112600301}$}
\fancyfoot[C]{\thepage\quad \textit{of}\quad\pageref*{LastPage}}

\fancypagestyle{plain}{
  \fancyhf{}
  \renewcommand{\headrulewidth}{0pt}
  \fancyfoot[C]{\thepage\quad \textit{of}\quad\pageref*{LastPage}}
}

\pgfplotsset{compat=1.18}
\usepackage{okumacro}

\renewcommand{\abstractname}{注意事項}

\newtagform{textbf}[\textbf]{[}{]}
\usetagform{textbf}

\newcommand*{\ie}{\textbf{\textit{i.e.}}\@\xspace}

\renewcommand{\qedsymbol}{$\blacksquare$}

\newtcbtheorem[]{reidai}{例題}
{fonttitle=\gtfamily\sffamily\bfseries\upshape\large,
 colframe=black,colback=black!15!white,
 rightrule=1pt,leftrule=1pt,bottomrule=2pt,
 colbacktitle=black,theorem style=standard,breakable,arc=10pt}
{tha}

\renewcommand{\thefootnote}{\arabic{footnote}}

\newtheoremstyle{mystyle}%
  {}{}{}{}{\bfseries}:{ }
  {\thmname{#1}\thmnumber{ #2}\thmnote{ (#3)}}

\theoremstyle{mystyle}

\newtheorem{dfn}{\texttt{Def.}}[section]
\newtheorem{exm}[dfn]{\texttt{Ex.}}
\newtheorem{prop}[dfn]{\texttt{Prop.}}
\newtheorem{lem}[dfn]{\texttt{Lem.}}
\newtheorem{thm}[dfn]{\texttt{Thm.}}
\newtheorem{cor}[dfn]{\texttt{Cor.}}
\newtheorem{rem}[dfn]{\texttt{Rem.}}
\newtheorem{fact}[dfn]{\texttt{Fact}}

\renewcommand{\qedsymbol}{$\blacksquare$}

\usepackage{lipsum}
\usepackage{float}

\newcommand{\kai}
{\noindent
\begin{tikzpicture}[scale=0.2, baseline=2.8pt]
\draw (3.3,1.2) node{\large\textgt{解 答}};
\draw[thick, rounded corners=3pt,] (0,0)--(6.5,0)--(6.5,2.4)--(0,2.4)--cycle;
\end{tikzpicture}}

\newcommand{\shomei}
{\noindent
\begin{tikzpicture}[scale=0.2, baseline=2.8pt]
\draw (3.3,1.2) node{\textgt{証 明}};
\draw[double,thick,rounded corners=3pt,] (0,0)--(6.5,0)--(6.5,2.4)--(0,2.4)--cycle;
\end{tikzpicture};}

\newcommand{\hosoku}{\noindent
\begin{tikzpicture}[scale=0.2, baseline=2.8pt]
\draw (6,1) node{\large\textgt{補足}};
\fill (0,1)--(1,0)--(2,1)--(1,2)--cycle;
\fill[gray] (1,1)--(2,0)--(3,1)--(2,2)--cycle;
\fill (2,1)--(3,0)--(4,1)--(3,2)--cycle;
\fill (10,1)--(11,0)--(12,1)--(11,2)--cycle;
\fill[gray] (9,1)--(10,0)--(11,1)--(10,2)--cycle;
\fill (8,1)--(9,0)--(10,1)--(9,2)--cycle;
\end{tikzpicture};}

\newcommand{\chui}{\noindent
\begin{tikzpicture}[scale=0.2, baseline=2.8pt]
\fill (0,0)--(6.5,0)--(6.5,2.2)--(0,2.2);
\draw (3.3,1) node[white]{\large\textgt{注意！}};
\draw[thick] (0,0)--(6.5,0)--(6.5,2.2)--(0,2.2)--cycle;
\end{tikzpicture};}

\title{\vspace{-3cm} \textbf{\texttt{REPORT 1(SPECIAL LECTURES)}}}
\author{\texttt{YI Ran} - $\mathnormal{61112600301}$\\ \texttt{andreyi@outlook.jp}}
\date{\today}

\begin{document}
\maketitle
\thispagestyle{plain}

%============================================================
%\section*{\textbf{\large Problem 1}\quad\normalfont\large}
%============================================================

% \begin{abstract}
% 本解答では，内積 $(\,\cdot\,|\,\cdot\,)$ は第2変数について線形，
% 第1変数について共役線形とする規約を採用する．
% また $H$ は $\dim H = n$ の有限次元複素 Hilbert 空間，
% $B(H)$ は $H$ 上の有界線形作用素全体を表す．
% 有限次元では任意の線形作用素が自動的に有界であるから，
% $B(H)$ は $H$ 上のすべての線形変換の集合と見なしてよい．
% \end{abstract}

%\vspace{0.5em}

%------------------------------------------------------------
\section*{\boxed{\textbf{\large 問1}}}
%------------------------------------------------------------

\textrm{Let $H$ be a Hilbert space with $\dim H = n$ and $M_n(\mathbb{C})$ be the vector space
of $n \times n$ matrices over $\mathbb{C}$.
$\mathrm{Tr}$ means the canonical trace on a matrix.
That is, for any $X = [x_{ij}] \in M_n(\mathbb{C})$,
$\mathrm{Tr}(X) = \sum_{i=1}^{n} x_{ii}$.}

\begin{enumerate}[(1)]
  \item \textrm{Let $\{e_i\}_{i=1}^{n}$ be an orthogonal basic system.
  That is, for any $x \in H$, $\exists\, \{x_i\}_{i=1}^{n} \subset \mathbb{C}$ such that
  $x = \sum_{i=1}^{n} x_i e_i$.
  Then show that the following map $\pi : B(H) \to M_n(\mathbb{C})$ is a linear isomorphism:
  \[
    \pi(T) = \bigl[(e_i \mid T e_j)\bigr]_{i,j=1}^{n}.
  \]}

  \item \textrm{Using the above map $\pi$ we define
  $\mathrm{Tr}(T) = \mathrm{Tr}(\pi(T)) = \sum_{i=1}^{n}(e_i \mid Te_i)$.
  Let $\{f_i\}_{i=1}^{n}$ be another orthogonal basic system for $H$.
  Show that $\mathrm{Tr}(T) = \sum_{i=1}^{n}(f_i \mid Tf_i)$.}
\end{enumerate}

\vspace{1em}

%------------------------------------------------------------
\section*{\boxed{\textbf{\large 解答}}}
%------------------------------------------------------------
\subsection*{\textbf{(1)}}
\begin{proof}
\ \\
\paragraph{\textbf{$\pi$の線形性}を示す}\ \\
\noindent
任意の $S, T \in B(H)$ と $\alpha, \beta \in \mathbb{C}$ をとる。$\pi$ の定義より、
$$
  \pi(\alpha S + \beta T) = \bigl[(e_i \mid (\alpha S + \beta T)e_j)\bigr]_{i,j=1}^{n}
$$
となる。ここで、各$i, j$ に対して，$\pi(\alpha S + \beta T)_{ij} = (e_i \mid (\alpha S + \beta T)e_j)$である。\\
\noindent
内積の線形性を用いると、
\[
    (e_i \mid (\alpha S + \beta T)e_j)
  = (e_i \mid \alpha S e_j + \beta T e_j)
  = \alpha (e_i \mid S e_j) + \beta (e_i \mid T e_j)
  = \bigl(\alpha \pi(S) + \beta \pi(T)\bigr)_{ij}
\]
これはすべての $i, j$ について成り立つので、
$\pi(\alpha S + \beta T) = \alpha \pi(S) + \beta \pi(T)$となり、$\pi$ は線形である。

\vspace{0.8em}
\noindent
\paragraph{\textbf{$\pi$ の単射性}を示す}\ \\
\noindent
$\forall S, T \in B(H)$を取り、$\pi(S)=\pi(T)$ と仮定する。このとき、
\[
  \pi(S-T)=\pi(S)-\pi(T)=0
\]
となる。\\
\noindent
したがって，すべての $i,j$ に対して、
\[
  (e_i\mid (S-T)e_j)=0
\]
が成り立つ。各 $j$ を固定し、$v = (S-T)e_j$ とする。$\forall i$に対して、
\[
  (e_i\mid (S-T)e_j) = (e_i\mid v) = 0
\]
を得る。$\{e_1,\dots,e_n\}$は正規直交基底であるから、

\[
  v=\sum_{j=1}^n x_j e_j
\]
ここで、
$$ 
  (e_i\mid v) = (e_i\mid \sum_{j=1}^n x_j e_j) = \sum_{j=1}^n x_j (e_i\mid e_j) =\sum_{j=1}^n x_j \delta_{ij} = x_i\qquad (\delta_{ij}=\begin{cases}1 & i=j,\\0 & i\ne j\end{cases})
$$
よって、
$$ 
  v = \sum_{j=1}^n x_j e_j = \sum_{j=1}^n (e_j\mid v) e_j = 0
$$
となる。次、$v = (S-T)e_j$を$ v=\sum_{j=1}^n (e_j\mid v) e_j$に代入すると、\\
$$
  (S-T)e_j = \sum_{j=1}^n (e_j\mid (S-T)e_j) e_j = 0
$$
となる。$(e_i\mid (S-T)e_j) = 0\quad \textrm{for }\ \forall i,j$より、
$$ 
  (S-T)e_j = \sum_{k=1}^n 0 \cdot e_j = 0
$$
となる。for $\forall x \in H$に対して、$x = \sum_{j=1}^n x_j e_j$と表せるから、
\[
  (S-T)x=\sum_{j=1}^n x_j (S-T)e_j= \sum_{j=1}^n x_j\cdot 0 =0
\]
したがって $S-T=0$、\ie $S=T$ である。ゆえに $\pi$ は単射である。

\vspace{0.8em}
\noindent
\paragraph{\textbf{$\pi$の全射性}を示す}\ \\
\noindent
任意の行列 $A = [a_{ij}] \in M_n(\mathbb{C})$ をとり、各基底ベクトル $e_j$ に対して、
\[
  T_A e_j := \sum_{i=1}^{n} a_{ij} e_i \qquad (j = 1, \ldots, n)
\]
とする。さらに、任意の $x=\sum_{j=1}^n x_je_j\in H$ に対して
\[
  T_Ax:=\sum_{j=1}^n x_j T_Ae_j
\]
とする。この定義より，任意の $e,h\in H$ および $\alpha,\beta\in\mathbb{C}$ に対して
\[
  T_A(\alpha e+\beta h)=\alpha T_Ae+\beta T_Ah
\]
が成り立つので，$T_A:H\to H$ は線形作用素である。
また，$\dim H=n$ であるから、$H$ は有限次元であり、
有限次元空間上の線形作用素は有界であるので $T_A\in B(H)$ である。次に、各 $k, j$ に対して、
\[
  \pi(T_A)_{kj}
  =(e_k \mid T_A e_j)
  = \left(e_k \;\middle|\; \sum_{i=1}^{n} a_{ij} e_i\right)
  = \sum_{i=1}^{n} a_{ij} (e_k \mid e_i)
  = \sum_{i=1}^{n} a_{ij} \delta_{ki}
  = a_{kj}
\]
したがって、すべての $k,j$ に対して $\pi(T_A)_{kj}=a_{kj}$ であるから、$\pi(T_A)=A$が成り立つ。よって $\pi$ は全射である。\\
\noindent
以上より、$\pi : B(H) \to M_n(\mathbb{C})$ は線形同型(Linear Isomorphism)である。
\end{proof}

\vspace{1.5em}

%------------------------------------------------------------
\subsection*{\textbf{(2)}}
%------------------------------------------------------------

\noindent
$\{f_i\}_{i=1}^{n}$ を $H$ の別の正規直交基底とし、$T\in B(H)$とする。
以下，$\mathrm{Tr}(T) = \sum_{i=1}^{n}(f_i \mid T f_i)$ を示す。\\

\begin{proof}\ \\
\noindent
$\{e_k\}_{k=1}^n$ は $H$ の正規直交基底であるから，各 $f_i$ は
\[
  f_i=\sum_{k=1}^n u_{ki}e_k\ ,
  \quad
  u_{ki}=(e_k\mid f_i)
\]
と表される。
ここで係数行列$U=(u_{ki})_{k,i=1}^n$を考える。\\
\noindent
$\{e_k\}_{k=1}^n$ と $\{f_i\}_{i=1}^n$ はいずれも正規直交基底であるから、
$U$ はユニタリ行列であり、
\[
  U^*U=I
\]
が成り立つ。
したがって、その $(k,l)$ 成分を比較すると
\[
  \sum_{i=1}^n \overline{u_{ki}}\,u_{li}=\delta_{kl}
  \qquad
  (k,l=1,\dots,n)
\]
を得る。

\medskip
\noindent
今示したいのは
\[
  \mathrm{Tr}(T)=\sum_{i=1}^n (f_i\mid Tf_i)
\]
である。
そこで右辺の和を計算すると、
\[
  \sum_{i=1}^n (f_i\mid Tf_i)
  =
  \sum_{i=1}^n
  \left(
    \sum_{k=1}^n u_{ki}e_k
    \;\middle|\;
    T\sum_{l=1}^n u_{li}e_l
  \right)
\]
となる。$T$ の線形性より
\[
  T\left(\sum_{l=1}^n u_{li}e_l\right)
  =
  \sum_{l=1}^n u_{li}Te_l
\]
であるから、
\[
  \sum_{i=1}^n (f_i\mid Tf_i)
  =
  \sum_{i=1}^n
  \left(
    \sum_{k=1}^n u_{ki}e_k
    \;\middle|\;
    \sum_{l=1}^n u_{li}Te_l
  \right).
\]
さらに、
\[
  \sum_{i=1}^n (f_i\mid Tf_i)
  =
  \sum_{i=1}^n \sum_{k=1}^n \sum_{l=1}^n
  \overline{u_{ki}}\,u_{li}\,(e_k\mid Te_l)
  =
  \sum_{k=1}^n \sum_{l=1}^n\bigl(\sum_{i=1}^n 
  \overline{u_{ki}}\,u_{li}\bigr)(e_k\mid Te_l)
\]

\medskip
$\sum_{i=1}^n \overline{u_{ki}}\,u_{li}=\delta_{kl}$を上式に代入すると
\[
  \sum_{i=1}^n (f_i\mid Tf_i)
  =
  \sum_{k=1}^n \sum_{l=1}^n \delta_{kl}(e_k\mid Te_l).
\]
となる。$\delta_{kl}$ は $k=l$ のときに $1$，$k\neq l$ のときに $0$ であるから、
右辺では $k=l$ の項だけが残り、
\[
  \sum_{i=1}^n (f_i\mid Tf_i)
  =
  \sum_{k=1}^n (e_k\mid Te_k).
\]
最後に，$\mathrm{Tr}(T)$ の定義より
\[
  \sum_{i=1}^n (f_i\mid Tf_i) = \sum_{k=1}^n (e_k\mid Te_k)=\mathrm{Tr}(\pi(T)) = \mathrm{Tr}(T)
\]
である。したがって、
\[
  \mathrm{Tr}(T)=\sum_{i=1}^n (f_i\mid Tf_i)
\]
が示された。
\end{proof}
\newpage
%------------------------------------------------------------
\section*{\boxed{\textbf{\large 問2}}}
%------------------------------------------------------------
\noindent
\textrm{Let $x_1, x_2, x_3 \in \mathbb{R}$ and let
\[
  A = \frac{1}{2}
  \begin{pmatrix}
    1 + x_3 & x_1 - ix_2 \\
    x_1 + ix_2 & 1 - x_3
  \end{pmatrix}.
\]
Show that $A$ is a density matrix (that is, $A \geq 0$ and $\mathrm{Tr}(A) = 1$)
if and only if $x_1^2 + x_2^2 + x_3^2 \leq 1$.}


%------------------------------------------------------------
\section*{\boxed{\textbf{\large 解答}}}
%------------------------------------------------------------


\begin{proof}\ \\

\paragraph{\textbf{$\mathrm{Tr}(A) = 1$ の確認}}\ \\
\noindent
トレースの定義より、
\[
  \mathrm{Tr}(A)
  = \frac{1}{2}(1 + x_3) + \frac{1}{2}(1 - x_3)
  = 1.
\]
よって $\mathrm{Tr}(A) = 1$ は $x_1, x_2, x_3$ の値によらず常に成り立つ。

\vspace{0.8em}
\noindent
\paragraph{\textbf{$A$ がエルミート行列であることの確認}}\ \\
\noindent
$x_1, x_2, x_3 \in \mathbb{R}$ より、共役転置を計算すると、
\[
  A^* = \overline{A^T} =\frac{1}{2}
  \begin{pmatrix}
    \overline{1 + x_3} & \overline{x_1 + ix_2} \\
    \overline{x_1 - ix_2} & \overline{1 - x_3}
  \end{pmatrix}
  = \frac{1}{2}
  \begin{pmatrix}
    1 + x_3 & x_1 - ix_2 \\
    x_1 + ix_2 & 1 - x_3
  \end{pmatrix}
  = A
\]
したがって $A = A^*$ であり、$A$ はエルミート行列である。

\vspace{0.8em}
\noindent
\paragraph{\textbf{Aの固有値の計算}}\ \\
\noindent
エルミート行列 $A$ が半正定値 ($A \geq 0$) であることは、
$A$ のすべての固有値が非負であることと同値である。
$A$ の特性多項式を計算する。
\[
  \det(\lambda I - A)
  = \left(\lambda - \frac{1+x_3}{2}\right)\!\left(\lambda - \frac{1-x_3}{2}\right)
    - \frac{x_1 - ix_2}{2} \cdot \frac{x_1 + ix_2}{2}.
\]

よって、
$$ 
  \det(\lambda I - A) = \lambda^2 - \lambda + \frac{1 - x_3^2}{4} - \frac{x_1^2 + x_2^2}{4} = \lambda^2 - \lambda + \frac{1 - (x_1^2 + x_2^2 + x_3^2)}{4}.
$$
$r = \sqrt{x_1^2 + x_2^2 + x_3^2} \geq 0$ とおくと、判別式は $1^2 - 4 \cdot 1 \cdot \dfrac{1-r^2}{4} = r^2 \geq 0$ となり、
固有値は実数であって、
\[
  \lambda_{\pm} =\frac{-(-1) \pm \sqrt{r^2}}{2 \cdot 1} = \frac{1 \pm |r|}{2} = \frac{1 \pm r}{2}
\]
となる。

\vspace{0.8em}
\noindent
\paragraph{$(\Rightarrow)$ \textbf{$A$ が密度行列} $\Rightarrow$ $x_1^2 + x_2^2 + x_3^2 \leq 1$:}\ \\
\noindent
$A$ が密度行列であると仮定する。すなわち $A \geq 0$ かつ $\mathrm{Tr}(A) = 1$ とする。
$A \geq 0$ より、すべての固有値は非負であるから、
\[
  \lambda_- = \frac{1 - r}{2} \geq 0\ , \quad \lambda_+ = \frac{1 + r}{2} \geq 0 
\]
これを解くと $0 \leq r \leq 1$、\ie $0 \leq r^2 \leq 1$であり、よって
\[
  x_1^2 + x_2^2 + x_3^2 \leq 1
\]
が従う。
\vspace{0.8em}
\noindent
\paragraph{$(\Leftarrow)$ $x_1^2 + x_2^2 + x_3^2 \leq 1$ $\Rightarrow$ \textbf{$A$ が密度行列}:}\ \\
\noindent
$x_1^2 + x_2^2 + x_3^2 \leq 1$、すなわち $0 \leq r=\sqrt{x_1^2 + x_2^2 + x_3^2} \leq 1$ と仮定する。
このとき、$0\leq r \leq 1$より、
$$
\lambda_+ = \dfrac{1 + r}{2} \geq 0\ ,\ \lambda_- = \dfrac{1 - r}{2} \geq 0
$$
が成り立つ。
よって、すべての固有値が非負であるから、 $A \geq 0$ が成り立つ。\\
\noindent
また、上で示したように $\mathrm{Tr}(A) = 1$ は常に成立する。
したがって $A$ は密度行列である。

\vspace{0.8em}
\noindent
以上 $(\Rightarrow)$ および $(\Leftarrow)$ より、$A$ が密度行列であること ($A \geq 0$ かつ $\mathrm{Tr}(A) = 1$) は
$x_1^2 + x_2^2 + x_3^2 \leq 1$ と同値である。\ie $A$が密度行列$\xLeftrightarrow{\text{iff}} x_1^2 + x_2^2 + x_3^2 \leq 1 $。
\end{proof}

\end{document}